%! Properties and Applications of Determinants — Unit 5, Lesson 5.2 — Exit Ticket (Answer Key)
\documentclass[10pt]{article}
\usepackage{linalg-article}
\usepackage{linalg-key}   % pulls in -boxes; do NOT also load -boxes
\newcommand{\TallMath}[1]{$\displaystyle #1\rule[-1.4em]{0pt}{3.2em}$}

\begin{document}

\pageheader{Unit 5, Lesson 5.2}{Exit Ticket --- Answer Key}
\namedateperiod

\vspace{0.15in}

No notes. Show your reasoning.

\begin{enumerate}[leftmargin=1.8em, itemsep=16pt]
  \item \textbf{Determinant by elimination.} Reduce to triangular and multiply the pivots to find
        $\det\begin{bmatrix} 1 & 2 & 0 \\ 3 & 7 & 1 \\ 0 & 2 & 5 \end{bmatrix}$. \; (No row swaps are needed.)
        \ansline{$R_2-3R_1\Rightarrow\begin{bsmallmatrix} 1 & 2 & 0\\ 0 & 1 & 1\\ 0 & 2 & 5\end{bsmallmatrix}$; $R_3-2R_2\Rightarrow\begin{bsmallmatrix} 1 & 2 & 0\\ 0 & 1 & 1\\ 0 & 0 & 3\end{bsmallmatrix}$. Pivots $1,1,3$, no swaps, so $\det=1\cdot1\cdot3=3$.}
  \item \textbf{Product rule.} A matrix $A$ has $\det A=3$ and a matrix $B$ has $\det B=5$. Find
        $\det(AB)$ and $\det A^{-1}$.
        \ansline{$\det(AB)=\det A\,\det B=3\cdot5=15$; \; $\det A^{-1}=1/\det A=1/3$.}
  \item \textbf{Explain.} You add $4\times(\text{row }1)$ to row $2$ of a matrix. What happens to its
        determinant, and \emph{why}? Use the word \emph{shear} in your one-sentence answer.
        \ansline{The determinant is \emph{unchanged}: adding a multiple of one row to another is a \emph{shear}, which slides the box without changing its base or height, so the volume (the determinant) stays the same.}
\end{enumerate}

\begin{teachernote}
Sort into ``got it / rule slip / sign slip.'' Item~1 is the fast method --- reduce to triangular and
multiply the pivots (watch for arithmetic in the elimination). Item~2 is the product rule and the
inverse consequence ($\det A^{-1}=1/\det A$). Item~3 is the target justification: shear $=$ no change,
because base and height are preserved. If many miss item~1, open 5.3 with a 60-second ``reduce and
multiply the pivots'' re-warm.
\end{teachernote}

\end{document}
